\documentclass[12pt,a4paper]{article}
\usepackage[utf8]{inputenc}
\usepackage[T1]{fontenc}
\usepackage{geometry}
\usepackage{graphicx}
\usepackage{float}
\usepackage{listings}
\usepackage{xcolor}
\usepackage{hyperref}
\usepackage{enumitem}
\usepackage{titlesec}
\usepackage{mdframed}
\usepackage{booktabs}
\geometry{margin=1in}

\titleformat{\section}
  {\normalfont\Large\bfseries}{\thesection.}{0.5em}{}
\titleformat{\subsection}
  {\normalfont\large\bfseries}{\thesubsection.}{0.5em}{}

\lstset{
  basicstyle=\ttfamily\small,
  breaklines=true,
  frame=single,
  backgroundcolor=\color{gray!10},
  keywordstyle=\color{blue},
  commentstyle=\color{green!60!black},
  stringstyle=\color{red},
  numbers=left,
  numberstyle=\tiny\color{gray},
  numbersep=8pt,
  showstringspaces=false,
  tabsize=2
}

\hypersetup{
  colorlinks=true,
  linkcolor=blue,
  urlcolor=blue
}

\begin{document}

%% ── Title Page ──────────────────────────────────────────────────────────────
\begin{titlepage}
\begin{center}
\vspace*{1cm}

{\Large \bfseries MEAN STACK \\[0.2cm] \large \normalfont(LPEIT-109)} \\[0.3cm]
{\large ASSIGNMENT}

\vspace{0.3cm}

{\large \bfseries GURU NANAK DEV ENGINEERING COLLEGE}

\vspace{0.5cm}

{\normalsize SUBMITTED IN THE PARTIAL FULFILLMENT OF THE REQUIREMENTS \\
FOR THE AWARD OF THE DEGREE OF \\[0.3cm]
{\large \bfseries BACHELORS OF TECHNOLOGY} \\[0.2cm]
(INFORMATION TECHNOLOGY)}}

\vspace{0.8cm}

\includegraphics[width=3in,height=3in]{assets/gnelogo.png}

\vspace{0.5cm}

{\large SUBMITTED BY: \\[0.3cm]
SHARNJEET SINGH (2303059)}

\vspace{1cm}

{\normalsize DEPARTMENT OF INFORMATION TECHNOLOGY \\[0.2cm]
GURU NANAK DEV ENGINEERING COLLEGE LUDHIANA \\[0.2cm]
(An Autonomous College Under UGC ACT)}
\end{center}
\end{titlepage}

\newpage

%% ── Q1 ───────────────────────────────────────────────────────────────────────
\section{Q1: Login Page Web Application Using MEAN Stack Authentication}

\subsection{Aim}
To construct a simple login page web application that authenticates users using the
complete MEAN stack — MongoDB, Express.js, Angular, and Node.js — with
\textbf{bcryptjs} for password hashing and \textbf{JSON Web Tokens (JWT)} for
stateless session management.

\subsection{Introduction}
Authentication is the process of verifying who a user is before granting access to
protected resources. In this project (\textit{The Elevated Navigator} — a Travel
Itinerary Planner), the login system is built across all four layers of the MEAN stack:

\begin{itemize}
  \item \textbf{MongoDB} stores user documents with bcrypt-hashed passwords.
  \item \textbf{Express.js} exposes \texttt{POST /api/auth/login} and
        \texttt{POST /api/auth/register} REST endpoints.
  \item \textbf{Angular} renders the login form, calls the API, stores the JWT,
        and guards protected routes.
  \item \textbf{Node.js} is the runtime that ties everything together.
\end{itemize}

\subsection{Packages Used}

\begin{center}
\begin{tabular}{lll}
\toprule
\textbf{Package} & \textbf{Layer} & \textbf{Purpose} \\
\midrule
bcryptjs        & Backend  & Hash passwords before saving to MongoDB \\
jsonwebtoken    & Backend  & Sign and verify JWT tokens \\
dotenv          & Backend  & Load \texttt{JWT\_SECRET} from \texttt{.env} \\
mongoose        & Backend  & User schema and database queries \\
@angular/common & Frontend & \texttt{HttpClient}, \texttt{isPlatformBrowser} \\
rxjs            & Frontend & Reactive \texttt{Observable} / \texttt{tap} operator \\
\bottomrule
\end{tabular}
\end{center}

\subsection{Backend — Auth Middleware (\texttt{server/middleware/auth.js})}

\begin{lstlisting}[language=JavaScript, caption={JWT helpers and protect middleware}]
const jwt     = require('jsonwebtoken');
const bcrypt  = require('bcryptjs');
const User    = require('../models/User');

const JWT_SECRET = process.env.JWT_SECRET || 'mean_mini_secret_key';

// Middleware: verify JWT on every protected route
const authenticate = async (req, res, next) => {
  try {
    const authHeader = req.headers['authorization'];
    // Header format: "Bearer <token>"
    const token = authHeader && authHeader.split(' ')[1];

    if (!token)
      return res.status(401).json({ message: 'Access token required' });

    // Decode and verify the token using the secret key
    const decoded = jwt.verify(token, JWT_SECRET);

    // Fetch the full user from DB (exclude password field)
    const user = await User.findById(decoded.id).select('-password');

    if (!user || !user.isActive)
      return res.status(401).json({ message: 'Invalid or expired token' });

    req.user = user;   // attach user to request for downstream use
    next();
  } catch (error) {
    return res.status(401).json({ message: 'Invalid or expired token' });
  }
};

// Sign a new JWT valid for 7 days
const generateToken = (userId) =>
  jwt.sign({ id: userId }, JWT_SECRET, { expiresIn: '7d' });

// Hash a plain-text password with bcrypt (salt rounds = 10)
const hashPassword = async (password) => {
  const salt = await bcrypt.genSalt(10);
  return bcrypt.hash(password, salt);
};

// Compare a plain-text password against a stored bcrypt hash
const comparePassword = async (plain, hashed) =>
  bcrypt.compare(plain, hashed);

module.exports = { authenticate, generateToken, hashPassword, comparePassword };
\end{lstlisting}

\subsection{Backend — Auth Routes (\texttt{server/routes/authRoutes.js})}

\begin{lstlisting}[language=JavaScript, caption={Register and Login endpoints}]
const express = require('express');
const User    = require('../models/User');
const { authenticate, generateToken,
        hashPassword, comparePassword } = require('../middleware/auth');

const router = express.Router();

// POST /api/auth/register
router.post('/register', async (req, res) => {
  try {
    const { name, email, password } = req.body;

    // Validate that all fields are present
    if (!name || !email || !password)
      return res.status(400).json({ success: false,
                                    message: 'All fields are required.' });

    const cleanEmail = String(email).trim().toLowerCase();

    // Prevent duplicate accounts
    const existing = await User.findOne({ email: cleanEmail });
    if (existing)
      return res.status(400).json({ success: false,
                                    message: 'User already exists.' });

    // Hash the password before storing — never save plain text
    const hashed = await hashPassword(password);

    const user = await User.create({
      name: String(name).trim(),
      email: cleanEmail,
      password: hashed,
      role: 'user',
    });

    // Issue a JWT immediately so the user is logged in after registering
    const token = generateToken(user._id);

    return res.status(201).json({
      success: true,
      token,
      user: { id: user._id, name: user.name,
              email: user.email, role: user.role },
    });
  } catch (error) {
    return res.status(500).json({ success: false,
                                  message: 'Server error during registration.' });
  }
});

// POST /api/auth/login
router.post('/login', async (req, res) => {
  try {
    const { email, password } = req.body;

    if (!email || !password)
      return res.status(400).json({ success: false,
                                    message: 'Email and password are required.' });

    // Look up the user by email
    const user = await User.findOne({
      email: String(email).trim().toLowerCase()
    });

    // Return the same generic error for both "not found" and "wrong password"
    // to avoid leaking which emails are registered
    if (!user || !user.isActive)
      return res.status(401).json({ success: false,
                                    message: 'Invalid email or password.' });

    // bcrypt.compare hashes the plain password and checks it against the stored hash
    const isMatch = await comparePassword(password, user.password);
    if (!isMatch)
      return res.status(401).json({ success: false,
                                    message: 'Invalid email or password.' });

    // Credentials are valid — sign and return a JWT
    const token = generateToken(user._id);

    return res.json({
      success: true,
      message: 'Login successful',
      token,
      user: { id: user._id, name: user.name,
              email: user.email, role: user.role },
    });
  } catch (error) {
    return res.status(500).json({ success: false,
                                  message: 'Server error during login.' });
  }
});

module.exports = router;
\end{lstlisting}

\subsection{Frontend — Auth Service (\texttt{frontend/src/app/services/auth.service.ts})}

\begin{lstlisting}[language=JavaScript, caption={Angular service — stores JWT and exposes reactive state}]
@Injectable({ providedIn: 'root' })
export class AuthService {
  private baseUrl   = 'http://localhost:5000/api/auth';
  // isPlatformBrowser prevents localStorage access during SSR prerendering
  private isBrowser = isPlatformBrowser(inject(PLATFORM_ID));

  // Angular signal — reactive current-user state used by navbar and guards
  currentUser = signal<AuthUser | null>(this.loadUser());

  // Read persisted user from localStorage on service init
  private loadUser(): AuthUser | null {
    if (!this.isBrowser) return null;
    try { return JSON.parse(localStorage.getItem('user') || 'null'); }
    catch { return null; }
  }

  get token()      { return this.isBrowser
                       ? localStorage.getItem('token') : null; }
  get isLoggedIn() { return !!this.token; }

  // Call the login endpoint; on success persist token + user
  login(email: string, password: string): Observable<any> {
    return this.http.post<any>(`${this.baseUrl}/login`,
                               { email, password }).pipe(
      tap(res => {
        if (res.success && this.isBrowser) {
          localStorage.setItem('token', res.token);
          localStorage.setItem('user', JSON.stringify(res.user));
          this.currentUser.set(res.user);   // update signal instantly
        }
      })
    );
  }

  // Clear token and redirect to login page
  logout(): void {
    if (this.isBrowser) {
      localStorage.removeItem('token');
      localStorage.removeItem('user');
    }
    this.currentUser.set(null);
    this.router.navigate(['/login']);
  }
}
\end{lstlisting}

\subsection{Frontend — HTTP Interceptor (\texttt{auth.interceptor.ts})}

\begin{lstlisting}[language=JavaScript, caption={Attach Bearer token to every outgoing request automatically}]
export const authInterceptor: HttpInterceptorFn = (req, next) => {
  const token = inject(AuthService).token;

  if (token) {
    // Clone the request and add the Authorization header
    const cloned = req.clone({
      setHeaders: { Authorization: `Bearer ${token}` },
    });
    return next(cloned);   // forward the modified request
  }

  return next(req);        // no token — send request as-is
};
\end{lstlisting}

\subsection{Frontend — Route Guard (\texttt{auth.guard.ts})}

\begin{lstlisting}[language=JavaScript, caption={Block unauthenticated access to /dashboard and /admin}]
export const authGuard: CanActivateFn = () => {
  const auth   = inject(AuthService);
  const router = inject(Router);

  if (auth.isLoggedIn) return true;   // token present — allow navigation

  router.navigate(['/login']);         // no token — redirect to login
  return false;
};
\end{lstlisting}

\subsection{Frontend — Login Component}

\begin{lstlisting}[language=JavaScript, caption={login.component.ts — calls AuthService on form submit}]
export class LoginComponent {
  email = ''; password = ''; error = ''; loading = false;

  constructor(private auth: AuthService, private router: Router) {}

  onSubmit() {
    if (!this.email || !this.password) return;
    this.loading = true;
    this.error   = '';

    this.auth.login(this.email, this.password).subscribe({
      next: (res) => {
        this.loading = false;
        // Navigate to dashboard on successful login
        if (res.success) this.router.navigate(['/dashboard']);
        else this.error = res.message || 'Login failed';
      },
      error: (err) => {
        this.loading = false;
        // Show the server error message (e.g. "Invalid email or password")
        this.error = err?.error?.message || 'Invalid credentials';
      }
    });
  }
}
\end{lstlisting}

\subsection{Authentication Flow Diagram}

\begin{mdframed}
\begin{center}
\texttt{
\begin{tabular}{l}
1. User fills email + password on Angular login form \\
2. Angular calls POST /api/auth/login via AuthService \\
3. Express finds user in MongoDB by email \\
4. bcrypt.compare(plain, hash) checks the password \\
5. On match: jwt.sign(\{id\}, SECRET, \{expiresIn:'7d'\}) \\
6. Token returned to Angular in JSON response \\
7. Angular stores token in localStorage \\
8. currentUser signal updated → navbar re-renders \\
9. HTTP Interceptor adds "Authorization: Bearer <token>" \\
   to every future request automatically \\
10. authGuard reads token to allow /dashboard access \\
\end{tabular}
}
\end{center}
\end{mdframed}

\subsection{Project Screenshots}

\textbf{Login Page:}
\begin{figure}[H]
  \centering
  \includegraphics[width=0.85\textwidth]{assets/screenshot-login.png}
  \caption{Login page of The Elevated Navigator}
\end{figure}

\textbf{Dashboard (after successful login):}
\begin{figure}[H]
  \centering
  \includegraphics[width=0.85\textwidth]{assets/screenshot-dashboard.png}
  \caption{User dashboard — accessible only after authentication}
\end{figure}

\textbf{Admin Panel:}
\begin{figure}[H]
  \centering
  \includegraphics[width=0.85\textwidth]{assets/screenshot-admin.png}
  \caption{Admin dashboard — role-based access control in action}
\end{figure}

\newpage

%% ── Q2 ───────────────────────────────────────────────────────────────────────
\section{Q2: Configure an Express Application with a Coding Example}

\subsection{Aim}
To configure an Express.js application, explaining each part of the setup with
in-code comments written in plain language.

\subsection{Introduction to Express}
Express is a minimal, unopinionated web framework for Node.js. It wraps Node's
built-in \texttt{http} module and adds:
\begin{itemize}
  \item A clean \textbf{routing API} (\texttt{app.get}, \texttt{app.post}, \ldots)
  \item A \textbf{middleware pipeline} — functions that run in sequence on every request
  \item Helpers for sending \textbf{JSON, HTML, files} and setting status codes
\end{itemize}

In this project the Express server lives in \texttt{server/server.js} and is the
backbone of the entire backend API.

\subsection{Step 1 — Install Dependencies}

\begin{lstlisting}[language=bash]
npm install express dotenv mongoose cors
\end{lstlisting}

\subsection{Step 2 — Environment File (\texttt{.env})}

\begin{lstlisting}[language=bash]
# Never commit this file — it holds secrets
MONGO_URI=mongodb://localhost:27017/travel_planner
PORT=5000
JWT_SECRET=your_super_secret_key
\end{lstlisting}

\subsection{Step 3 — Database Connection (\texttt{server/config/db.js})}

\begin{lstlisting}[language=JavaScript, caption={Mongoose connection helper}]
const mongoose = require('mongoose');

// Async function so we can use await and catch errors cleanly
const connectDB = async () => {
  const mongoUri = process.env.MONGO_URI;

  // Stop the process immediately if the URI is missing from .env
  if (!mongoUri) {
    console.error('MONGO_URI is missing in .env');
    process.exit(1);
  }

  try {
    // mongoose.connect returns a promise — await it before continuing
    await mongoose.connect(mongoUri, {
      useNewUrlParser:    true,   // use the new URL string parser
      useUnifiedTopology: true,   // use the new server discovery engine
    });
    console.log('MongoDB connected');
  } catch (error) {
    // If the DB is unreachable, log the reason and exit
    console.error('MongoDB connection error:', error.message);
    process.exit(1);
  }
};

module.exports = connectDB;
\end{lstlisting}

\subsection{Step 4 — Main Server File (\texttt{server/server.js})}

\begin{lstlisting}[language=JavaScript, caption={Fully commented Express configuration}]
// Load environment variables from .env into process.env
// Must be called before any other code reads process.env
const dotenv = require('dotenv');
dotenv.config();

// Core Node.js module for working with file paths
const path    = require('path');

// Express — the web framework
const express = require('express');

// Our custom DB connection helper
const connectDB = require('./config/db');

// Route handlers — each file exports an express.Router()
const authRoutes        = require('./routes/authRoutes');
const itineraryRoutes   = require('./routes/itineraryRoutes');
const userRoutes        = require('./routes/userRoutes');
const roleRequestRoutes = require('./routes/roleRequestRoutes');

// ── Create the Express application instance ──────────────────────────────────
const app = express();

// ── Connect to MongoDB ────────────────────────────────────────────────────────
// This runs once at startup; if it fails the process exits (see db.js)
connectDB();

// ── Global Middleware ─────────────────────────────────────────────────────────

// Parse incoming requests with a JSON body (e.g. from Angular HttpClient)
// Without this, req.body would be undefined for POST/PUT requests
app.use(express.json());

// Parse URL-encoded bodies (e.g. traditional HTML form submissions)
app.use(express.urlencoded({ extended: false }));

// Serve everything inside the /views folder as static files
// e.g. GET /login.html will return views/login.html directly
app.use(express.static(path.join(__dirname, '..', 'views')));

// ── Page Routes (serve HTML files) ───────────────────────────────────────────

// Root URL — send the landing page
app.get('/', (req, res) => {
  res.sendFile(path.join(__dirname, '..', 'views', 'index.html'));
});

// Login page
app.get('/login', (req, res) => {
  res.sendFile(path.join(__dirname, '..', 'views', 'login.html'));
});

// Register page
app.get('/register', (req, res) => {
  res.sendFile(path.join(__dirname, '..', 'views', 'register.html'));
});

// Dashboard page (Angular SPA handles auth on the client side)
app.get('/dashboard', (req, res) => {
  res.sendFile(path.join(__dirname, '..', 'views', 'dashboard.html'));
});

// ── API Routes ────────────────────────────────────────────────────────────────
// app.use(prefix, router) mounts a router at a URL prefix.
// All routes defined inside authRoutes will be prefixed with /api/auth
// e.g. router.post('/login') becomes POST /api/auth/login

app.use('/api/auth',          authRoutes);        // register, login, profile
app.use('/api/itinerary',     itineraryRoutes);   // CRUD + booking
app.use('/api/users',         userRoutes);        // user management
app.use('/api/role-requests', roleRequestRoutes); // admin role requests

// ── Start the Server ──────────────────────────────────────────────────────────
const PORT = process.env.PORT || 5000;

// Special case: Vercel runs the app as a serverless function,
// so we export the app instead of calling listen()
if (process.env.VERCEL) {
  module.exports = app;
} else {
  // app.listen binds to the port and starts accepting connections
  app.listen(PORT, () => {
    console.log(`Server running on port ${PORT}`);
  });
}
\end{lstlisting}

\subsection{Step 5 — A Sample API Route (\texttt{routes/itineraryRoutes.js})}

\begin{lstlisting}[language=JavaScript, caption={Express Router with middleware chaining}]
const express    = require('express');
const Itinerary  = require('../models/Itinerary');
const { authenticate, authorize } = require('../middleware/auth');

// express.Router() creates a mini-app that handles a subset of routes
const router = express.Router();

// GET /api/itinerary
// authenticate runs first — if the token is invalid it sends 401 and stops here
// If the token is valid, control passes to the async handler
router.get('/', authenticate, async (req, res) => {
  try {
    // Regular users only see active itineraries; admins see everything
    const query = req.user.role === 'user' ? { isActive: true } : {};

    // .populate replaces the createdBy ObjectId with the actual User document
    const itineraries = await Itinerary.find(query)
      .populate('createdBy', 'name email role')
      .sort({ createdAt: -1 });   // newest first

    return res.json(itineraries);
  } catch (error) {
    return res.status(500).json({ message: 'Server error.' });
  }
});

// POST /api/itinerary
// Two middleware functions before the handler:
//   1. authenticate — verify JWT
//   2. authorize    — check role is admin or superadmin
router.post('/', authenticate, authorize('admin', 'superadmin'),
  async (req, res) => {
    try {
      const { title, destination, startDate,
              endDate, duration, budget } = req.body;

      // Basic validation before touching the database
      if (!title || !destination || !startDate ||
          !endDate || !duration || budget === undefined)
        return res.status(400).json({ message: 'Fill all required fields.' });

      // Create the document; createdBy is taken from the verified JWT user
      const itinerary = await Itinerary.create({
        ...req.body,
        createdBy: req.user._id,
      });

      return res.status(201).json(itinerary);
    } catch (error) {
      return res.status(500).json({ message: 'Server error.' });
    }
  }
);

// Export the router so server.js can mount it
module.exports = router;
\end{lstlisting}

\subsection{How Middleware Chaining Works}

When a request arrives at \texttt{POST /api/itinerary}, Express runs each function
in the chain left-to-right:

\begin{mdframed}
\begin{center}
\texttt{
\begin{tabular}{l}
Request arrives \\
  $\downarrow$ \\
express.json()         -- parse the JSON body \\
  $\downarrow$ \\
authenticate()         -- verify the JWT token \\
  $\downarrow$ \\
authorize('admin',...)  -- check the user's role \\
  $\downarrow$ \\
Route handler          -- create the itinerary in MongoDB \\
  $\downarrow$ \\
res.json()             -- send the response back \\
\end{tabular}
}
\end{center}
\end{mdframed}

Each middleware calls \texttt{next()} to pass control to the next function, or calls
\texttt{res.status(...).json(...)} to short-circuit the chain and send an error
response immediately.

\subsection{Key Express Concepts Summary}

\begin{center}
\begin{tabular}{ll}
\toprule
\textbf{Concept} & \textbf{What it does in this project} \\
\midrule
\texttt{express()}            & Creates the application instance \\
\texttt{app.use()}            & Registers global middleware or mounts a router \\
\texttt{express.json()}       & Parses \texttt{application/json} request bodies \\
\texttt{express.static()}     & Serves the \texttt{views/} folder as static files \\
\texttt{express.Router()}     & Groups related routes into a separate file \\
\texttt{app.listen(PORT)}     & Starts the HTTP server on the given port \\
\texttt{req.body}             & Parsed request payload (needs \texttt{express.json()}) \\
\texttt{req.params.id}        & URL parameter, e.g. \texttt{/itinerary/:id} \\
\texttt{req.user}             & Custom property set by \texttt{authenticate} middleware \\
\texttt{res.json()}           & Sends a JSON response with \texttt{Content-Type: application/json} \\
\texttt{res.status(404)}      & Sets the HTTP status code before sending \\
\bottomrule
\end{tabular}
\end{center}

\newpage

%% ── GitHub Link ──────────────────────────────────────────────────────────────
\section*{GitHub Repository}

The complete source code for this project (The Elevated Navigator — MEAN Stack
Travel Itinerary Planner) is available at:

\begin{center}
  \large
  \href{https://github.com/sharnjeet21/mean_mini}{https://github.com/sharnjeet21/mean\_mini}
\end{center}

\vspace{0.5cm}
\noindent The repository contains:
\begin{itemize}
  \item \texttt{server/} — Express + Node.js backend with JWT auth and Mongoose models
  \item \texttt{frontend/} — Angular 21 SPA with Tailwind CSS design system
  \item \texttt{views/} — Legacy static HTML pages
  \item \texttt{.env.example} — Environment variable template
  \item \texttt{README.md} — Setup and run instructions
\end{itemize}

\end{document}
